\chapter{Linear continua and structural elements}\label{ch:elem}

\begin{tcolorbox}[colback=gray!5,colframe=gray!50,title={In this chapter}]
A thin steel strip is meshed with a standard plate element, loaded at the tip, and solved. The answer comes out two hundred times too stiff, and no setting is wrong: the element is simply too stiff for thin plates under the integration rule that looks the most careful. This chapter starts from that failure. It first shows how a weak form becomes a stiffness and a mass matrix, and which element families \fea{} offers for scalar fields, plane and solid continua, beams, plates and shells. It does not re-derive shape functions. It then looks at the places where an element can be wrong while looking plausible, namely locking, integration rules and the choice of mass matrix, and tests each against an independent reference: patch tests on distorted meshes, pure bending, a slender cantilever of a nearly incompressible material, beam frequencies against closed forms, and a plate against the shell that must reproduce it.\par\smallskip
After it you should be able to say how a weak form becomes a stiffness and a mass matrix, recognise locking or a wrong integration rule from a test with an independent reference, and choose between a consistent and a lumped mass matrix, knowing the frequency error of each.
\end{tcolorbox}

Imagine a designer who wants the tip deflection of a thin plate strip, $1\,\mathrm{m}$ long and $1\,\mathrm{cm}$ thick, clamped at one end. The finite element answer is $0.005$ of the beam-theory value on a coarse mesh, and still $0.6$ on a mesh sixteen times finer. Nothing is broken, and every element is positive definite and passes the patch test. The cure is a change of integration rule on one term of the element matrix, and the rest of the chapter explains why that works and how to tell it from the outside. The thin cantilever strip is the running example of this chapter, and its numbers appear again in \Cref{sec:el-locking}.

\section{From the weak form to an element matrix}\label{sec:el-weak}

The smallest complete example is the scalar element \anchor{fea_engine.elements.scalar.Quad4Laplace}, which has one unknown per node. For the equation $-\nabla\!\cdot(\kappa\nabla u)=q$ on a domain $\Omega$ with $u=0$ on a part of its boundary, the weak form is: find $u$ such that for every admissible $v$
\begin{equation}\label{eq:el-weak}
\int_\Omega\kappa\,\nabla v\cdot\nabla u\,d\Omega=\int_\Omega v\,q\,d\Omega .
\end{equation}
Inserting the interpolation $u\approx\sum_a N_a u_a$ and the same functions for $v$ turns \eqref{eq:el-weak} into $\bm K\bm u=\bm f$ with
\begin{equation}\label{eq:el-K}
\bm K_e=\int_{\Omega_e}\bm B^{T}\kappa\,\bm B\,d\Omega,\qquad \bm B=\nabla\bm N,\qquad
\bm f_e=\int_{\Omega_e}\bm N^{T}q\,d\Omega .
\end{equation}
The integrals are computed by Gauss quadrature on the reference element, with $d\Omega=\det\bm J\,d\xi$ and $\bm B$ mapped by the Jacobian $\bm J$. The elasticity elements have the same structure with a vector field: $\bm K_e=\int\bm B^{T}\bm D\bm B\,d\Omega$, where $\bm B$ is the strain-displacement matrix and $\bm D$ the elasticity matrix \citep{zienkiewicz2013fem,bathe2014fep}. Everything else in this chapter is a choice of $\bm N$, of $\bm D$, or of the quadrature rule in \eqref{eq:el-K}.

\begin{sidenote}{Strong form, weak form and virtual work}
The \emph{strong form} is the differential equation itself, required to hold at every point; it needs the second derivative of $u$. The \emph{weak form} multiplies the equation by a test function $v$, integrates over the domain and moves one derivative from $u$ onto $v$ by integration by parts, so that only first derivatives remain. Functions that are continuous across element boundaries but have a kink there, which is what polynomial pieces glued at the nodes are, can then be used. In mechanics $v$ is a virtual displacement, and \eqref{eq:el-weak} says that the internal virtual work equals the external virtual work for every admissible virtual displacement. Choosing $v=N_a$ for each node $a$ in turn gives one balance equation per degree of freedom (dof), the number of independent unknowns of the model.
\end{sidenote}

\tikzfigure[0.98\linewidth]{ch03_weak2matrix}{From the weak form to the matrices. The choices that decide whether an element is right are the interpolation, the material matrix and the quadrature rule.}{fig:el-weak2matrix}

\begin{algbox}{Element stiffness and mass by Gauss quadrature}
\begin{pseudo}
\Require node coordinates $\bm x_e$, material matrix $\bm D$ (or $\kappa$), density $\rho$, Gauss rule $\{\bm\xi_g,w_g\}$
\State $\bm K_e\gets\bm 0$;\ $\bm M_e\gets\bm 0$
\For{each Gauss point $g$}
  \State $\bm N,\ \partial\bm N/\partial\bm\xi\gets$ shape functions at $\bm\xi_g$
  \State $\bm J\gets(\partial\bm N/\partial\bm\xi)\,\bm x_e$;\ $\det\bm J$;\ $\bm B\gets\bm J^{-1}\partial\bm N/\partial\bm\xi$ mapped into strain--displacement form
  \State $\bm K_e\gets\bm K_e+w_g\det\bm J\;\bm B^T\bm D\bm B$
  \State $\bm M_e\gets\bm M_e+w_g\det\bm J\;\rho\bm N^T\bm N$
\EndFor
\If{lumped mass requested}
  \State $\bm M_e\gets$ HRZ lumping of $\bm M_e$ \Comment{diagonal, same total mass}
\EndIf
\Ensure $\bm K_e$, $\bm M_e$
\end{pseudo}
\end{algbox}

The rule evaluates $n_g$ integrands of $O(d^2n_{str})$ each, and it sets the rank of $\bm K_e$. It goes wrong when it is too low (hourglass modes, rank deficiency) or too high for the physics (locking, \Cref{sec:el-locking}), and when the Jacobian is not positive. The code is \anchor{fea_engine.elements.base.Element} (\code{stiffness}, \code{mass}) and \code{assemble_lumped_mass}; the checks are \Cref{sec:el-verif} (patch tests, beams) and \Cref{fig:el-mass} (mass choices).

\paragraph{An element in five steps.} The algorithm above packs five ideas into a few lines. (1)~\emph{Interpolation}: the displacement inside an element is built from its nodal values, $u(\bm x)=\sum_aN_a(\bm x)u_a$, where each shape function $N_a$ equals 1 at node $a$ and 0 at the others, so the unknowns are the nodal values and nothing else. (2)~\emph{Isoparametric map}: the same functions interpolate the geometry, $\bm x=\sum_aN_a\bm x_a$, so every element is computed on one reference element (for example $\xi\in[-1,1]$) and mapped to its place. The Jacobian $\bm J=\partial\bm x/\partial\bm\xi$ is the local stretch of that map, and $\det\bm J$ turns a reference volume into a physical one; it must be positive, which is what Chapter~\ref{ch:mesh} checks. (3)~\emph{$\bm B$ and $\bm D$}: the strain-displacement matrix $\bm B$ turns the nodal values into strain at a point, $\bm\varepsilon=\bm B\bm u_e$, and the material matrix $\bm D$ turns strain into stress, $\bm\sigma=\bm D\bm\varepsilon$. $\bm B$ carries the geometry and the interpolation, $\bm D$ the material. (4)~\emph{Quadrature}: the integral over the element is replaced by a weighted sum over Gauss points, which gives $\bm K_e=\sum_gw_g\det\bm J\,\bm B^T\bm D\bm B$. (5)~\emph{Assembly}: the element matrices are added into the global matrix at the rows and columns of the element's dofs, and the elements interact only through the nodes they share.

The smallest example is a two-node bar of length $L$, cross-section $A$ and modulus $E$. Its shape functions are $N_1=(1-\xi)/2$ and $N_2=(1+\xi)/2$, so $\bm J=L/2$, $\bm B=[-1/L,\ 1/L]$ and $\bm K_e=\frac{EA}{L}\left[\begin{smallmatrix}1&-1\\-1&1\end{smallmatrix}\right]$: stretching the bar by $\delta$ needs the force $EA\delta/L$, the familiar spring. The consistent mass matrix is $\bm M_e=\frac{\rho AL}{6}\left[\begin{smallmatrix}2&1\\1&2\end{smallmatrix}\right]$, and the lumped one puts half of the mass, $\rho AL/2$, on each node. \Cref{lst:el-bar} subclasses \anchor{fea_engine.elements.base.Element} with just the shape functions and $\bm B$ (the construction of Chapter~\ref{ch:arch}) and lets the package build the three matrices by quadrature.

\begin{lstlisting}[caption={The five steps for a two-node bar: the package builds $\bm K_e$, $\bm M_e$ and the lumped mass from the shape functions and $\bm B$ alone (run with the package on the path).},label={lst:el-bar}]
import numpy as np
from fea_engine.elements import Element, jacobian

class Bar(Element):                       # 2 nodes, one displacement per node
    n_nodes, dofs_per_node, dim, gauss_order = 2, 1, 1, 2
    dof_names = ("u",)
    def shape_and_derivs(self, xi):
        return np.array([(1 - xi[0]) / 2, (1 + xi[0]) / 2]), np.array([[-0.5, 0.5]])
    def B_matrix(self, xi, x):
        _, dN = self.shape_and_derivs(xi)
        J, detJ = jacobian(dN, x)
        return np.linalg.solve(J, dN), detJ      # B = dN/dx = J^-1 dN/dxi

E_, A, rho, L = 210e9, 1.0e-4, 7800.0, 2.0
x = np.array([[0.0], [L]])                # the element's two node coordinates
bar = Bar()
Ke = bar.stiffness(x, np.array([[E_ * A]]))
Me = bar.mass(x, np.array([[rho * A]]))
Ml = bar.lumped_mass(x, np.array([[rho * A]]))

np.set_printoptions(precision=4, suppress=True)
print("Ke / (EA/L):\n", Ke / (E_ * A / L))
print("Me / (rho A L / 6):\n", Me / (rho * A * L / 6))
print("lumped / (rho A L / 2):\n", Ml / (rho * A * L / 2))
print("total mass:", Me.sum(), Ml.sum(), "exact", rho * A * L)
\end{lstlisting}

\begin{lstlisting}[numbers=none,language={},keywordstyle={},commentstyle={},frame=single,aboveskip=2pt]
Ke / (EA/L):
 [[ 1. -1.]
 [-1.  1.]]
Me / (rho A L / 6):
 [[2. 1.]
 [1. 2.]]
lumped / (rho A L / 2):
 [[1. 0.]
 [0. 1.]]
total mass: 1.56 1.56 exact 1.56
\end{lstlisting}

The printed matrices are the hand-derived ones, and the lumped matrix keeps the total mass, $\rho AL=1.56$\,kg for these values. Two Gauss points are used because the mass integrand $N_aN_b$ is quadratic; the stiffness integrand is constant, so one point would be enough for $\bm K_e$.

\begin{sidenote}{Admissible function and quadrature}
The function $v$ in the weak form is a \emph{test function}: any function that is zero where $u$ is prescribed. Taking $v=N_a$ for each node $a$ in turn gives one equation per node. The integrals in \eqref{eq:el-K} are rarely done exactly; a Gauss rule with $n_g$ points per direction integrates polynomials up to degree $2n_g-1$ exactly, and which rule is used is a modelling choice, not just a numerical one, as the next sections show.
\end{sidenote}

\paragraph{The mass matrix: consistent or lumped.}
In words: the consistent matrix spreads the mass of the element continuously, so that the nodes are coupled (the off-diagonal terms $\rho AL/6$ of the bar above), while lumping replaces the element by point masses at its nodes, which decouples them. The same interpolation gives the consistent mass matrix $\bm M_e=\int\rho\bm N^{T}\bm N\,d\Omega$, which is sparse with the pattern of $\bm K$. A lumped mass matrix is diagonal: \fea{} provides element-by-element HRZ lumping (\code{assemble_lumped_mass}), which keeps the total mass of the element. The diagonal form is what the explicit time integration of Chapter~\ref{ch:dyn} needs, since it makes the solve for the accelerations trivial. The choice has a measurable effect on the eigenvalues. For the lowest eigenvalue of the unit-square membrane, $\lambda_1=2\pi^2$ for a clamped edge, \Cref{fig:el-mass} shows that the consistent mass overestimates it and the lumped mass underestimates it, both with the error falling like $h^{\cecMassConsOrder}$ and $h^{\cecMassLumpOrder}$. On the finest mesh the consistent mass is high by \cecMassConsFine{} and the lumped mass is low by \cecMassLumpFine{}. The signs are the usual ones for these two choices, and an average of the two is a classical way to reduce the error; this chapter does not test that.

\begin{sidenote}{Natural frequency and the eigenvalue problem}
A structure vibrating freely without damping moves as $\bm u(t)=\bm\phi\,e^{\mathrm i\omega t}$, which turns $\bm M\ddot{\bm u}+\bm K\bm u=\bm0$ into the generalized eigenproblem $\bm K\bm\phi=\lambda\bm M\bm\phi$ with $\lambda=\omega^2$. An eigenvalue $\lambda$ gives a natural frequency $f=\sqrt\lambda/2\pi$ in hertz, and its eigenvector $\bm\phi$ is the mode shape, the pattern in which the structure moves at that frequency. Only the lowest few modes are of interest in this book, and the same problem appears in the eigenvalues of the membrane below, where $\bm K$ and $\bm M$ come from the scalar element. Chapter~\ref{ch:dyn} treats it in full.
\end{sidenote}

\begin{figure}[ht]
\centering
\includegraphics[width=0.55\linewidth]{figures/ch03_mass.pdf}
\caption{Lowest eigenvalue of the clamped unit-square membrane (\code{Quad4Laplace}), consistent against lumped mass: the relative error in absolute value. The consistent-mass eigenvalue is above the exact one and the lumped-mass eigenvalue is below it.}\label{fig:el-mass}
\end{figure}

\section{The element families}\label{sec:el-families}

\Cref{tab:el-elements} lists the elements treated in this chapter with their nodes, degrees of freedom per node and the \code{gauss_order} attribute of the class. Plane and solid elements share one idea: a displacement-based isoparametric element with polynomial interpolation. The plane elements use $\bm D$ for plane stress or plane strain; the 3-D elements use the $6\times6$ isotropic $\bm D$. Linear triangles and tetrahedra have constant strain; the 4-node quadrilateral and the 8-node hexahedron have a bilinear or trilinear field and therefore represent bending badly (\Cref{sec:el-locking}).

\begin{sidenote}{Plane stress, plane strain and the matrix $\bm D$}
The plane elements reduce a three-dimensional body to two dimensions in one of two ways. \emph{Plane stress} describes a thin body loaded in its plane, such as a plate: the stress out of the plane is zero and the thickness is free to change. \emph{Plane strain} describes a long body held in the third direction, such as a dam section: the strain out of the plane is zero and the stress there is not. In both cases the in-plane strains are $(\varepsilon_x,\varepsilon_y,\gamma_{xy})$ and $\bm D$ is a $3\times3$ matrix, for plane stress $\frac{E}{1-\nu^2}\left[\begin{smallmatrix}1&\nu&0\\\nu&1&0\\0&0&\frac{1-\nu}2\end{smallmatrix}\right]$, for plane strain $\frac{E}{(1+\nu)(1-2\nu)}\left[\begin{smallmatrix}1-\nu&\nu&0\\\nu&1-\nu&0\\0&0&\frac{1-2\nu}2\end{smallmatrix}\right]$. The plane-strain matrix grows without bound as the Poisson ratio $\nu$ approaches $0.5$, which is the root of the volumetric locking of \Cref{sec:el-locking}.
\end{sidenote}

\begin{table}[ht]
\centering\small
\begin{tabular}{llrrr}
\toprule
family & element & nodes & dofs per node & \code{gauss\_order} \\
\midrule
scalar & \code{Quad4Laplace} & 4 & 1 & 2 \\
plane & \code{Tri3PlaneStress} & 3 & 2 & 1 \\
plane & \code{Tri6PlaneStress} & 6 & 2 & 2 \\
plane & \code{Quad4PlaneStress} & 4 & 2 & 2 \\
plane & \code{Quad8PlaneStress} & 8 & 2 & 3 \\
solid & \code{Tet4Solid3D} & 4 & 3 & 1 \\
solid & \code{Tet10Solid3D} & 10 & 3 & 2 \\
solid & \code{Hex8Solid3D} & 8 & 3 & 2 \\
solid & \code{Hex8SolidBbar} & 8 & 3 & 2 \\
solid & \code{Hex20Solid3D} & 20 & 3 & 3 \\
beam & \code{Beam2DEulerBernoulli} & 2 & 2 & 1 \\
beam & \code{Beam2DReissner} & 2 & 3 & 2 \\
beam & \code{Beam3DEulerBernoulli} & 2 & 6 & 2 \\
plate/shell & \code{Quad4MindlinPlate} & 4 & 3 & 2 \\
plate/shell & \code{Shell4MITC} & 4 & 6 & 2 \\
\bottomrule
\end{tabular}

\caption{Elements of this chapter. \code{gauss_order} is the attribute of the class as it stands in tag \booktag; its meaning is the number of points per direction for quadrilaterals and hexahedra, and a simplex rule order for triangles and tetrahedra.}\label{tab:el-elements}
\end{table}

\paragraph{Structural elements.}
The beam elements use a one-dimensional mesh. \anchor{fea_engine.elements.beams.Beam2DEulerBernoulli} has two dofs per node (transverse displacement and rotation) and cubic Hermite interpolation. \anchor{fea_engine.elements.beams.Beam2DReissner} adds shear deformation, so it has an axial displacement as well, and \anchor{fea_engine.elements.beams3d.Beam3DEulerBernoulli} has six dofs per node. The corotational versions of these beams belong to Chapter~\ref{ch:nonlin}. \anchor{fea_engine.elements.plates.Quad4MindlinPlate} is a 4-node plate with deflection and two rotations per node, and \anchor{fea_engine.elements.shells.Shell4MITC} is a 4-node shell with membrane, bending and transverse-shear action and six dofs per node, with the assumed transverse shear strain of \citet{dvorkin1984}.

\paragraph{Which element for which structure.} \Cref{tab:el-which} says what each family assumes about the structure. The assumption is the reason for the element's unknowns, and it also says when the element is the wrong choice.

\begin{table}[ht]
\centering\small
\begin{tabularx}{\linewidth}{@{}>{\raggedright\arraybackslash}p{0.2\linewidth} >{\raggedright\arraybackslash}p{0.33\linewidth} >{\raggedright\arraybackslash}X@{}}
\toprule
family & what it assumes & use it for \\
\midrule
plane or solid continuum & nothing about the shape of the deformation; the full strain state is computed & compact bodies, details, and any region where a beam or plate theory does not hold \\
Euler--Bernoulli beam & cross-sections stay plane and perpendicular to the axis, so there is no shear deformation; unknowns are the deflection and the rotation & slender beams, where the depth is much smaller than the length \\
Reissner (Timoshenko) beam & sections stay plane but may rotate relative to the axis, so shear deformation is included & thicker beams, and higher modes whose wavelength approaches the depth \\
Mindlin plate & lines normal to the mid-surface stay straight but not normal; unknowns are the deflection and two rotations & thin to moderately thick flat plates \\
shell & a curved surface carrying membrane, bending and transverse-shear action together & curved thin-walled structures \\
\bottomrule
\end{tabularx}
\caption{What the structural element families assume.}\label{tab:el-which}
\end{table}

\begin{sidenote}{Beam and plate theories in one paragraph}
The classical (Euler--Bernoulli) beam and the Kirchhoff plate assume that a line normal to the neutral axis or mid-surface remains normal after bending, which makes the shear strain zero and the rotation equal to the slope of the deflection. A Kirchhoff plate element therefore needs the slope to be continuous between elements, which is hard to arrange. The Timoshenko--Reissner beam and the Reissner--Mindlin plate drop the normality assumption, keep the line straight, and treat the rotation as an independent unknown; only continuity of the unknowns, not of their slopes, is then needed, so the elements are simple. The price is that, as the structure gets thin, the shear strain must tend to zero, and an element that cannot reproduce that limit locks (\Cref{sec:el-locking}). The MITC family of \citet{dvorkin1984} avoids this by interpolating the transverse shear strains from values at chosen points of the element edges (``tying points''), instead of computing them at Gauss points from the displacement field; the name stands for mixed interpolation of tensorial components. Textbooks that treat all of this are \citet{zienkiewicz2013fem}, \citet{bathe2014fep} and \citet{cook2002}.
\end{sidenote}

\section{Locking and how the elements avoid it}\label{sec:el-locking}

Locking is an element behaving too stiffly for a deformation mode that its interpolation cannot reproduce without also producing a parasitic strain. It is the one error in this chapter that grows when the problem gets \emph{easier} (a thinner beam, a more incompressible material), which is why it needs a test of its own.

\begin{sidenote}{What locking looks like inside an element}
Take a 4-node quadrilateral in pure bending. The bending field needs the top and bottom edges to curve, but the edges of a bilinear element stay straight, so the element can only approximate bending by shearing, and the shear strain it develops is parasitic: no shear exists in pure bending. That strain stores energy, the element resists bending more than the real structure does, and the result is a stiffness that is too high. Refining the mesh reduces the effect slowly. Locking is therefore not a bug in the code but a property of the formulation, and the cures all change the formulation: a different integration rule, a different interpolation of the troublesome strain, or a higher order.
\end{sidenote}

\Cref{tab:el-intrules} lists the three integration choices that appear in the rest of this section: full integration, which locks; one-point integration, which has hourglass modes; and selective reduced integration, which is the cure used by the plate element below.

\begin{table}[ht]
\centering\small
\begin{tabularx}{\linewidth}{@{}>{\raggedright\arraybackslash}p{0.2\linewidth} >{\raggedright\arraybackslash}p{0.16\linewidth} >{\raggedright\arraybackslash}p{0.25\linewidth} >{\raggedright\arraybackslash}X@{}}
\toprule
integration rule & Gauss points per direction & rank of $\bm K_e$ (4-node quadrilateral) & usual symptom \\
\midrule
full & 2 & 5: the three rigid-body modes are the only zero-energy ones & locks in bending and for nearly incompressible material \\
reduced & 1 & 3: two extra zero-energy modes & hourglass modes: mesh deformations that cost no energy and grow unchecked \\
selective reduced & 2 for the energy terms that do not lock, 1 for the one that does & 5 & usually neither; the choice of the terms is part of the element \\
\bottomrule
\end{tabularx}
\caption{The three integration choices for a bilinear element. The rank is the number of independent strain states the element energy can see; the quadrilateral has 8 dofs and 3 rigid-body modes, so the correct rank is 5. Hourglass control (stabilisation) is described by \citet{belytschko2014nonlinear}.}\label{tab:el-intrules}
\end{table}

\paragraph{Try it.} \Cref{lst:el-try} meshes the strip of the introduction with \code{Quad4MindlinPlate} and solves it twice on each mesh: with full integration for both the bending and the shear energy (the careful-looking choice) and with the package default, selective reduced integration. The ratio is the tip deflection divided by the beam-theory value.

\begin{lstlisting}[caption={A thin cantilever strip with two integration rules for the same plate element (run with the package on the path).},label={lst:el-try}]
from fea_engine import FESystem, Material, Quad4MindlinPlate
from fea_engine.material import D_mindlin_plate
from fea_engine.mesh import rectangle_mesh

mat = Material(E=2.1e11, nu=0.3, rho=7850.0)
h, L, W, P = 0.01, 1.0, 0.2, -1000.0                  # thin strip: h/L = 0.01
Db, Ds = D_mindlin_plate(mat, h)
ref = P * L**3 / (3 * mat.E * h**3 / (12 * (1 - mat.nu**2)) * W)   # beam theory with plate rigidity

for nx in (4, 16, 64):
    mesh = rectangle_mesh(L, W, nx, max(2, nx // 4))
    left, tip = mesh.nodes_on_plane(axis=0, value=0.0), mesh.nodes_on_plane(axis=0, value=L)
    out = {}
    for rule in ("full", "sri"):                      # same element, two integration rules
        s = FESystem(mesh, Quad4MindlinPlate())
        s.assemble_stiffness((Db, Ds), integration=rule)
        s.fix_dofs(left, [0, 1, 2]); s.add_nodal_force(tip, 0, P)
        out[rule] = s.solve_static()[3 * tip[0]] / ref
    print(f"{nx:3d} elements: tip deflection / beam theory   "
          f"full {out['full']:.3f}   selective-reduced {out['sri']:.3f}")
\end{lstlisting}

\begin{lstlisting}[numbers=none,language={},keywordstyle={},commentstyle={},frame=single,aboveskip=2pt]
  4 elements: tip deflection / beam theory   full 0.005   selective-reduced 1.059
 16 elements: tip deflection / beam theory   full 0.081   selective-reduced 1.077
 64 elements: tip deflection / beam theory   full 0.602   selective-reduced 1.078
\end{lstlisting}

The fully integrated element is two hundred times too stiff on four elements and still $40\,\%$ too stiff on sixty-four, so refining the mesh is an expensive and poor remedy. The selective rule is within $6\,\%$ at once. The reference is approximate, which is why the converged ratio is $1.078$ and not $1$; the difference between the two rows is what matters. The mechanism, the same one in continuum elements, is the next topic.

\tikzfigure[0.98\linewidth]{ch03_locking}{The mechanism of shear locking in a bilinear element under bending (schematic). The exact bending mode curves the edges and has no shear strain; the bilinear interpolation keeps the edges straight, which produces a parasitic shear strain and an artificially stiff element.}{fig:el-lockmech}

\paragraph{Bending of continuum elements.}
A 4-node quadrilateral in bending cannot represent the curvature of a beam without shear strain, and the parasitic shear stiffens the beam. \Cref{fig:el-bending} loads a plane-stress beam by a pure end moment and plots the ratio of the computed tip deflection to the exact one (the exact solution is quadratic in the displacements, so it lies in the space of the quadratic elements). The linear elements converge from below: at 4 elements along the beam the ratio is \cecBendQuadFour{} for \code{Quad4PlaneStress} and \cecBendTriFour{} for \code{Tri3PlaneStress}, and at 32 elements it is \cecBendQuadThirtytwo{} and \cecBendTriThirtytwo{}. The observed order of the quadrilateral error is \cecBendQuadOrder, close to the expected second order. The quadratic elements \code{Quad8PlaneStress} and \code{Tri6PlaneStress} reproduce the exact solution up to rounding (\cecBendQuadEightErr{} and \cecBendTriSixErr{} for the tested meshes), as the theory says they must.

\begin{figure}[ht]
\centering
\includegraphics[width=0.6\linewidth]{figures/ch03_bending.pdf}
\caption{Pure bending of a plane-stress beam, tip deflection divided by the exact value, against the number of elements along the beam. The quadratic elements are exact on these meshes. The helper is the one used by \texttt{fea\_engine/tests/test\_benchmarks\_convergence.py}.}\label{fig:el-bending}
\end{figure}

\begin{sidenote}{Incompressibility and the Poisson ratio}
The Poisson ratio $\nu$ is the ratio of lateral shortening to axial stretching in a bar. For $\nu=0.5$ the volume does not change under load, which is nearly true of rubber ($\nu\approx0.499$) and of metals in plastic flow. The bulk modulus $K=E/\bigl(3(1-2\nu)\bigr)$ measures the resistance to volume change and tends to infinity as $\nu\to0.5$, so any volume change costs an energy $\tfrac12K\varepsilon_v^2$ that dwarfs everything else. A low-order element must then keep the volumetric strain $\varepsilon_v=\operatorname{tr}\bm\varepsilon$ close to zero at every Gauss point; it has more such constraints than free dofs to satisfy them with, and it can only comply by not deforming.
\end{sidenote}

\paragraph{Volumetric locking and mean-dilatation B-bar.}
The idea of the cure is to constrain the volume change once per element, not once per Gauss point. A nearly incompressible material ($\nu\to0.5$) forces $\operatorname{tr}\bm\varepsilon\to0$ at every Gauss point, which a low-order element satisfies only by not deforming. The mean-dilatation (B-bar) element of \citet{hughes1980bbar}, \anchor{fea_engine.elements.solids.Hex8SolidBbar}, replaces the volumetric part of $\bm B$ by its average over the element, which relaxes the constraint to one condition per element. \Cref{tab:el-locking} bends a slender cantilever ($L/h=20$, square section) with a tip load and compares the tip deflection with the Euler--Bernoulli value. At $\nu=0.3$ the plain 8-node hexahedron gives \cecLockHexPlainA{} of the beam value (shear locking, as in two dimensions), and when $\nu\to0.4999$ it collapses to \cecLockHexPlainB{}, because volumetric locking is added. The B-bar element gives \cecLockHexBbarA{} at $\nu=0.3$ and \cecLockHexBbarB{} at $\nu=0.4999$: it removes the \emph{additional} locking caused by incompressibility and leaves the shear locking untouched, which is the claim of the method and nothing more. The 20-node hexahedron gives \cecLockHexTwentyA{} at $\nu=0.3$ and \cecLockHexTwentyB{} at $\nu=0.4999$ on a mesh with the same node spacing, so it locks little in shear but still shows some volumetric locking.

\begin{table}[ht]
\centering\small
\begin{tabular}{lrrrr}
\toprule
element & mesh & dofs & $\nu=0.3$ & $\nu=0.4999$ \\
\midrule
\code{Hex8Solid3D} & $20\times 4\times 4$ & 1575 & 0.711 & 0.177 \\
\code{Hex8SolidBbar} & $20\times 4\times 4$ & 1575 & 0.729 & 0.745 \\
\code{Hex20Solid3D} & $10\times 2\times 2$ & 963 & 0.989 & 0.905 \\
\bottomrule
\end{tabular}

\caption{Tip deflection of a slender cantilever divided by the Euler--Bernoulli value, three hexahedra and two Poisson ratios. The 20-node mesh has the node spacing of the 8-node mesh. A 3-D beam does not reach exactly 1 even for a perfect element (root and shear effects), so only the differences between the rows are significant.}\label{tab:el-locking}
\end{table}

\paragraph{Shear locking in the plate and its cure.}
\anchor{fea_engine.elements.plates.Quad4MindlinPlate} uses full $2\times2$ Gauss integration for the bending energy and a reduced one-point rule for the transverse shear energy (selective reduced integration, the default \code{'sri'}). With full integration for both, the element locks as the plate gets thin: \Cref{tab:el-plate} gives the tip deflection of a cantilever strip ($h/L=0.01$, $W/L=0.2$) divided by the beam-theory value $PL^3/(3D_pW)$ with the plate rigidity $D_p=Eh^3/(12(1-\nu^2))$. This reference is approximate (it ignores the anticlastic effect and shear), so the converged ratio is \cecPlateSriFine{} and not 1; what the table shows is the difference between the columns. With selective reduced integration the ratio is \cecPlateSriCoarse{} on a mesh of 4 elements along the length. With full integration it is \cecPlateFullCoarse{} on the same mesh and still \cecPlateFullFine{} on 64 elements. The shell \code{Shell4MITC}, which uses a different cure for locking (assumed shear strains tied at the element edges), converges to the same value: on the finest mesh the ratio of the shell to the plate is \cecPlateRatioFS. A flat shell must reproduce the plate, so this agreement is a check of both elements; the membrane response of the shell in this test was exactly zero, as it must be for a flat homogeneous single layer.

\begin{table}[ht]
\centering\small
\begin{tabular}{rrrr}
\toprule
elements along $L$ & \code{Quad4MindlinPlate} (SRI) & same, full integration & \code{Shell4MITC} \\
\midrule
4 & 1.059 & 0.005 & 1.051 \\
8 & 1.081 & 0.021 & 1.069 \\
16 & 1.077 & 0.081 & 1.075 \\
32 & 1.078 & 0.262 & 1.077 \\
64 & 1.078 & 0.602 & 1.078 \\
\bottomrule
\end{tabular}

\caption{Tip deflection of a flat cantilever strip divided by the beam-theory value with plate rigidity (an approximate reference), against the number of elements along the length (a quarter as many across).}\label{tab:el-plate}
\end{table}

\section{Verification}\label{sec:el-verif}

\begin{sidenote}{Convergence order and mesh refinement}
Refining the mesh means making the element size $h$ smaller. For a smooth problem the error behaves like $Ch^r$, so a plot of the error against $h$ on logarithmic axes is a straight line whose slope is the \emph{order of convergence} $r$, the number the tables and figures of this chapter report as the observed order. For elements that contain complete polynomials of degree $p$ the displacement error is of order $p+1$ (in the $L^2$ norm) and the error in the strains and stresses is of order $p$. For the eigenvalues of a problem of order $2m$ ($m=1$ for a membrane or a bar, $m=2$ for a beam) it is of order $2(p+1-m)$ \citep{strangfix1973}, so the lowest frequency of the cubic Hermite beam ($p=3$, $m=2$) converges with order 4. A measured order that is lower than the expected one is a warning: either the solution is not smooth, or the element has a defect.
\end{sidenote}

\begin{sidenote}{Patch test}
A patch test asks whether an element can reproduce the simplest nonzero strain state, constant strain, exactly on an irregular patch of elements. An element that cannot is not guaranteed to converge to the right answer when the mesh is refined, so the test is a necessary condition. It says nothing about locking, since a locking element still passes it.
\end{sidenote}

\paragraph{Patch tests.}
An element passes the patch test if a mesh of distorted elements reproduces a linear displacement field exactly when the boundary nodes are given the field, and recovers the constant stress. \Cref{tab:el-patch} gives the largest relative error of the displacement and of the stress for the \cecPatchCount{} continuum elements on distorted meshes; the worst value is \cecPatchWorst{}, which is rounding. The patch test is necessary but not sufficient for convergence, and it says nothing about locking, which is why the tests above exist. The B-bar hexahedron is not in this table: it does not support stress recovery in the nodal form the test uses, so its patch test (displacements only) is the one of \texttt{tests/test\_hex8\_bbar.py}, which passes together with the other three tests of that file (4 passed in the run for this chapter).

\begin{table}[ht]
\centering\small
\begin{tabular}{lrr}
\toprule
element & displacement error & stress error \\
\midrule
\code{Quad4PlaneStress} & $5.5\times10^{-16}$ & $5.2\times10^{-15}$ \\
\code{Quad8PlaneStress} & $1.0\times10^{-15}$ & $7.6\times10^{-15}$ \\
\code{Tri3PlaneStress} & $3.7\times10^{-16}$ & $4.9\times10^{-15}$ \\
\code{Tri6PlaneStress} & $1.8\times10^{-15}$ & $1.7\times10^{-14}$ \\
\code{Hex8Solid3D} & $1.8\times10^{-16}$ & $8.9\times10^{-16}$ \\
\code{Hex20Solid3D} & $1.2\times10^{-15}$ & $2.7\times10^{-15}$ \\
\code{Tet4Solid3D} & $3.6\times10^{-16}$ & $8.9\times10^{-16}$ \\
\code{Tet10Solid3D} & $1.2\times10^{-15}$ & $1.9\times10^{-15}$ \\
\bottomrule
\end{tabular}

\caption{Patch tests on distorted meshes. The test helper is \texttt{\_patch} of \texttt{test\_benchmarks\_convergence.py}.}\label{tab:el-patch}
\end{table}

\paragraph{Beams.}
The Euler--Bernoulli element is exact for an end load: one element reproduces the tip deflection $PL^3/(3EI)$ to \cecEBTipErr, because the cubic Hermite interpolation contains the exact solution. For vibration the cubic Hermite element converges with order \cecEBFreqOrder{} in the lowest cantilever frequency (\cecEBFreqHz{} Hz for the test beam), the expected fourth order. The Reissner beam approaches the Timoshenko tip deflection $PL^3/(3EI)+PL/(GA)$ with order \cecTimoOrder{} (\Cref{fig:el-beams}); shear contributes \cecTimoShearShare{}\% of the tip deflection for this slender test beam, so it is a check of the formulation more than a demonstration of shear effects.

\begin{figure}[ht]
\centering
\includegraphics[width=0.95\linewidth]{figures/ch03_beams.pdf}
\caption{Convergence of the first cantilever frequency of the Euler--Bernoulli beam (left) and of the tip deflection of the Reissner beam (right) with the element size, with the observed order.}\label{fig:el-beams}
\end{figure}

\begin{table}[ht]
\centering\small
\begin{tabularx}{\linewidth}{@{}>{\raggedright\arraybackslash}p{0.44\linewidth} X >{\centering\arraybackslash}p{0.11\linewidth}@{}}
\toprule
Concept & Code & Script calls \\
\midrule
scalar element, weak form, mass & \anchor{fea_engine.elements.scalar.Quad4Laplace} & yes \\
plane and solid elements & \anchor{fea_engine.elements.solids.Quad4PlaneStress}, \anchor{fea_engine.elements.solids.Quad8PlaneStress}, \anchor{fea_engine.elements.solids.Tri3PlaneStress}, \anchor{fea_engine.elements.solids.Tri6PlaneStress}, \anchor{fea_engine.elements.solids.Hex8Solid3D}, \anchor{fea_engine.elements.solids.Hex20Solid3D}, \anchor{fea_engine.elements.solids.Tet4Solid3D}, \anchor{fea_engine.elements.solids.Tet10Solid3D} & yes \\
B-bar hexahedron & \anchor{fea_engine.elements.solids.Hex8SolidBbar} & yes \\
beams & \anchor{fea_engine.elements.beams.Beam2DEulerBernoulli}, \anchor{fea_engine.elements.beams.Beam2DReissner}, \anchor{fea_engine.elements.beams3d.Beam3DEulerBernoulli} & yes \\
plate, shell & \anchor{fea_engine.elements.plates.Quad4MindlinPlate}, \anchor{fea_engine.elements.shells.Shell4MITC} & yes \\
convergence studies & \anchor{fea_engine.convergence.run_study}, \anchor{fea_engine.convergence.observed_order} & yes \\
\bottomrule
\end{tabularx}
\caption{Where the concepts of this chapter live in the repository (tag \booktag). The last column says whether this chapter's own checks call, by name, the code listed in the row: yes (every name), part, or no. Code reached only through another routine counts as no, so the column is a lower bound on what this book exercises.}\label{tab:el-map}
\end{table}

\begin{limitsbox}
\begin{itemize}
\item \textbf{Locking is not removed by the patch test.} All elements of \Cref{tab:el-patch} pass it, and the plain 8-node hexahedron still gives \cecLockHexPlainA{} of the beam value. For slender or nearly incompressible problems, choose the element by the tests of \Cref{sec:el-locking}, not by the patch test.
\item \textbf{B-bar addresses volumetric locking only.} It does not cure shear locking (\Cref{tab:el-locking}), and the 20-node element still showed some volumetric locking at $\nu=0.4999$ on the mesh used.
\item \textbf{Stress recovery} is not available for every element (for example \code{Hex8SolidBbar}); use \code{at='centroid'} where the error message says so.
\item \textbf{The plate reference is approximate.} \Cref{tab:el-plate} compares with a strip of plate rigidity, so only differences between its columns are claimed.
\item \textbf{Not covered:} orthotropic and layered materials, thermal and piezoelectric coupling, shells in large rotation (Chapter~\ref{ch:nonlin} treats the corotational and the director shell only in outline), and meshes of high aspect ratio or strong distortion beyond the patch-test meshes.
\end{itemize}
\end{limitsbox}

\section{Summary}\label{sec:el-summary}

We began with a thin strip that came out two hundred times too stiff. The cause is that every element is the integral $\int\bm B^{T}\bm D\bm B$ evaluated by Gauss quadrature, and the quadrature rule is part of the model: a rule that looks the most careful produced locking, and a cheaper rule applied to the shear term only removed it. The same quadrature gives the mass matrix, consistent or lumped, and the consistent mass overestimates the membrane eigenvalue while the lumped mass underestimates it.

All continuum elements pass the patch test on distorted meshes to rounding, yet the low-order ones lock in bending, and a nearly incompressible material locks the plain hexahedron to \cecLockHexPlainB{} of the beam value. B-bar fixes the incompressible part only, selective reduced integration keeps the Mindlin plate usable when thin, and the MITC shell agrees with the plate in the flat limit. The beams are exact (Euler--Bernoulli under an end load) or converge at the expected orders. The next chapter takes these matrices into time.

\section*{Exercises}\label{sec:el-exercises}
\begin{enumerate}
\item \emph{(Implementation)} Repeat \Cref{fig:el-mass} with the average of the consistent and the lumped matrix, $\bm M=\tfrac12(\bm M_c+\bm M_l)$. Does the error of $\lambda_1$ cancel to leading order, and what order does it have?
\item \emph{(Implementation)} Run the locking test of \Cref{tab:el-locking} on three meshes, $10\times2\times2$, $20\times4\times4$ and $40\times8\times8$, for the plain and the B-bar hexahedron. Does shear locking go away with refinement, and at what rate?
\item \emph{(Implementation)} Use \code{integration='full'} for the plate on a thicker strip ($h/L=0.1$). Is the locking still visible? Compare with the Timoshenko value.
\item \emph{(Interpretation)} Replace the end load of the cantilever by a distributed load. Is one Euler--Bernoulli element still exact? Why not?
\item \emph{(Implementation)} Distort the patch-test mesh more strongly (move the interior nodes by half an element). Does any element in \Cref{tab:el-patch} fail?
\item \emph{(Derivation)} For a fixed--free bar of one two-node element, find the lowest eigenvalue with the consistent mass matrix $\tfrac{\rho AL}{6}\left[\begin{smallmatrix}2&1\\1&2\end{smallmatrix}\right]$ and with the lumped matrix, and show that their errors against the exact $(\pi/2)^2EA/(\rho AL^2)$ are about $+21.6\,\%$ and $-18.9\,\%$.
\end{enumerate}

\chapterrefs
